\section*{Capítulo \ref{ch:b3:rate-theories} --- Teorías de la velocidad de reacción}

\begin{solution}{exo:b3:rate-theories:1}
$\bar v_{\mathrm{rel}} = \sqrt{8 \times 1.381\times10^{-23} \times 298/(\pi \times 18.46 \times 1.661\times10^{-27})} = \qty{585}{m/s}$.
$\sigma\bar v_{\mathrm{rel}}N_A = 0.43\times10^{-18} \times 585 \times 6.022\times10^{23} = \qty{1.51e8}{m^3.mol^{-1}.s^{-1}} =
\qty{1.5e11}{L.mol^{-1}.s^{-1}}$. Valor medido: $2.07\times10^{-12} \times 6.022\times10^{23}\times10^{-3} =
\qty{1.25e9}{L.mol^{-1}.s^{-1}}$, luego $P = 0.008$: menos de una colisión de cada cien tiene la
orientación adecuada (el O del NO debe encontrar un O terminal del \ce{O3}).
\end{solution}

\begin{solution}{exo:b3:rate-theories:2}
$\Delta^{\ddagger}H^\circ = 75.0 - 2.48 = \qty{72.5}{kJ/mol}$ en disolución; $75.0 - 4.96 = \qty{70.0}{kJ/mol}$
para una reacción bimolecular en fase gaseosa.
\end{solution}

\begin{solution}{exo:b3:rate-theories:3}
Diels--Alder: muy negativa (dos moléculas forman un solo complejo cíclico).
Pérdida de \ce{N2}: positiva (un enlace se afloja y se gana libertad). Creación de iones en un
disolvente polar: negativa, las cargas que se desarrollan ordenan el disolvente a su alrededor.
\end{solution}

\begin{solution}{exo:b3:rate-theories:4}
Agua: $8 \times 8.314 \times 298.15/(3 \times \num{0.890e-3}) = \qty{7.4e9}{L.mol^{-1}.s^{-1}}$; hexano:
\qty{2.2e10}{L.mol^{-1}.s^{-1}}. El hexano, tres veces menos viscoso: las moléculas difunden tres
veces más deprisa.
\end{solution}

\begin{solution}{exo:b3:rate-theories:5}
$\ln(k_2T_1/k_1T_2) = \ln(10 \times 298.15/318.15) = 2.238$; $\Delta^{\ddagger}H^\circ = 8.3145 \times 2.238/(1/298.15 -
1/318.15) = \qty{88.2}{kJ/mol}$. $\Delta^{\ddagger}S^\circ = R[\ln(kh/k_BT) + \Delta^{\ddagger}H^\circ/RT] = 8.3145 \times
(-36.365 + 35.596) = \qty{-6.4}{J.K^{-1}.mol^{-1}}$ (a \qty{298.15}{K}). $E_a = R\ln10/(1/298.15 - 1/318.15) =
\qty{90.8}{kJ/mol} \approx \Delta^{\ddagger}H^\circ + RT$ a la temperatura media.
\end{solution}

\begin{solution}{exo:b3:rate-theories:6}
$\mu_{\mathrm{OH}} = 16/17$, $\mu_{\mathrm{OD}} = 32/18$, razón de números de onda $\sqrt{1.889} = 1.374$:
$\tilde\nu_{\mathrm{OD}} = \qty{2661}{cm^{-1}}$. $k_{\mathrm H}/k_{\mathrm D} = \exp(1.4388 \times 996/(2 \times 298)) = \eu^{2.40} = 11$.
\end{solution}

\begin{solution}{exo:b3:rate-theories:7}
Máximo: $\exp(1.4388 \times 808/(2 \times 298)) = \eu^{1.95} = 7.0$. Un valor de 40 no puede proceder de las
\omterm{def:b3:quantum-model-systems:oscillator}{energías de punto cero}: el hidrógeno atraviesa la barrera por \omterm{def:b3:quantum-model-systems:tunnelling}{efecto túnel}, cosa que el deuterio,
el doble de pesado, hace mucho menos.
\end{solution}

\begin{solution}{exo:b3:rate-theories:8}
$\ce{S2O8^{2-}}$ e \ce{I-}: $z_{\mathrm A}z_{\mathrm B} = +2$, $\log(k/k_0) = 2 \times 0.51 \times 2 \times 0.10 = 0.20$: $k$
se multiplica por 1.6. Éster y \ce{OH-}: un compañero neutro, sin efecto.
$\ce{[Co(NH3)5Br]^{2+}}$ y \ce{OH-}: $z_{\mathrm A}z_{\mathrm B} = -2$, $k$ se multiplica por $10^{-0.20} = 0.63$.
\end{solution}

\begin{solution}{exo:b3:rate-theories:9}
Temprana, en el valle de entrada (un estado de transición parecido a los reactivos, según el postulado
de Hammond). La reacción inversa, endotérmica, tiene una barrera tardía en su propio
camino: la energía de vibración en el enlace H--F es la que mejor la favorece.
\end{solution}

\begin{solution}{exo:b3:rate-theories:10}
$\Delta^{\ddagger}H^\circ = -Rb = \qty{61.9}{kJ/mol}$, con un error típico de $R \times 31 = \qty{0.26}{kJ/mol}$ e
intervalo del 95\,\% de $\pm\qty{0.8}{kJ/mol}$.
\[
  \Delta^{\ddagger}S^\circ = R(16.528 - 23.760) = \qty{-60.1}{J.K^{-1}.mol^{-1}},
\]
con un error típico de $R \times 0.103 = \qty{0.86}{J.K^{-1}.mol^{-1}}$ e intervalo del 95\,\% de $\pm\qty{2.7}{J.K^{-1}.mol^{-1}}$.
\end{solution}

\begin{solution}{exo:b3:rate-theories:11}
A partir de $\ln k = \ln(k_BT/h) + \Delta^{\ddagger}S^\circ/R - \Delta^{\ddagger}H^\circ/RT$,
\[
  \frac{\dd\ln k}{\dd T} = \frac1T + \frac{\Delta^{\ddagger}H^\circ}{RT^2}, \qquad
  E_a = RT^2\frac{\dd\ln k}{\dd T} = \Delta^{\ddagger}H^\circ + RT .
\] Entonces $A = k\eu^{E_a/RT} = (k_BT/h)\eu^{\Delta^{\ddagger}S^\circ/R}\eu^{-\Delta^{\ddagger}H^\circ/RT}\eu^{(\Delta^{\ddagger}H^\circ + RT)/RT} =
\eu\,(k_BT/h)\eu^{\Delta^{\ddagger}S^\circ/R}$. Con $\Delta^{\ddagger}S^\circ = 0$: $A = 2.718 \times \num{6.21e12} = \qty{1.7e13}{s^{-1}}$, el
orden de los factores observados; los valores superiores indican una $\Delta^{\ddagger}S^\circ$ positiva, un
estado de transición laxo.
\end{solution}

\begin{solution}{exo:b3:rate-theories:12}
Con funciones de partición moleculares por unidad de volumen: $q^{\mathrm T}/V = (2\pi mk_BT)^{3/2}/h^3$ para
A, B y el complejo (masa $m_{\mathrm A} + m_{\mathrm B}$), y para el complejo $q^{\mathrm R} = 8\pi^2Ik_BT/h^2$
($\sigma = 1$). La razón de traslación es $h^3/(2\pi\mu k_BT)^{3/2}$, luego
\[
  k = \frac{k_BT}{h}\cdot\frac{h^3}{(2\pi\mu k_BT)^{3/2}}\cdot\frac{8\pi^2\mu d^2k_BT}{h^2}\,\eu^{-E_0/RT}
    = \frac{8\pi^2d^2(k_BT)^{1/2}}{(2\pi)^{3/2}\mu^{1/2}}\eu^{-E_0/RT}
    = \pi d^2\sqrt{\frac{8k_BT}{\pi\mu}}\,\eu^{-E_0/RT},
\]
por par de moléculas; multiplicar por $N_A$ da el \cref{thm:b3:rate-theories:collision}
con $\sigma = \pi d^2$.
\end{solution}

\begin{solution}{pb:b3:rate-theories:1}
\textbf{1.} $\mu_{\mathrm{CH}} = 12 \times 1.00783/13.00783 = \qty{0.9297}{u}$; $\mu_{\mathrm{CD}} = 12 \times 2.01410/14.01410 =
\qty{1.7246}{u}$.
\textbf{2.} $2917/\sqrt{1.7246/0.9297} = 2917/1.3620 = \qty{2142}{cm^{-1}}$, un 1.5\,\% por encima del
valor medido, 2109 (anarmonicidad y acoplamiento de los cuatro enlaces de \ce{CD4}).
\textbf{3.} $\frac12 \times 2917 = \qty{1458.5}{cm^{-1}}$ y $\frac12 \times 2109 = \qty{1054.5}{cm^{-1}}$; diferencia:
\qty{404}{cm^{-1}} $= \qty{4.83}{kJ/mol}$.
\textbf{4.} Se convierte en la coordenada de reacción: su \omterm{def:b3:quantum-model-systems:oscillator}{energía de punto cero} desaparece.
\textbf{5.} La barrera es más alta para D en \qty{4.83}{kJ/mol}.
\textbf{6.} $\exp(1.43878 \times 404/310.15) = \eu^{1.874} = 6.5$.
\textbf{7.} $\exp(1.43878 \times 404/298.15) = 7.0$.
\textbf{8.} En un estado de transición real, parte de la \omterm{def:b3:quantum-model-systems:oscillator}{energía de punto cero} de la tensión
sobrevive en otros modos, lo que reduce el efecto; las \omterm{def:b3:quantum-model-systems:oscillator}{energías de punto cero} por sí solas
no pueden superar este valor (el \omterm{def:b3:quantum-model-systems:tunnelling}{efecto túnel} sí).
\textbf{9.} Poco: son efectos secundarios, del orden del 10\,\% por deuterio como
mucho.
\textbf{10.} Que el enlace C--H se rompe en la etapa determinante de la velocidad, con la
tensión totalmente convertida en la coordenada de reacción.
\textbf{11.} Eliminación de primer orden: $t_{1/2} = \ln2/k_{\mathrm{total}}$, inversamente proporcional a ella.
\textbf{12.} $k_{\mathrm{total}}(\mathrm H) = 0.75\,k_{\mathrm{total}}(\mathrm H) + 0.25\,k_{\mathrm{total}}(\mathrm H)$: desmetilación y otras vías.
\textbf{13.} Se divide por 6.5.
\textbf{14.} $0.25 + 0.75/6.52 = 0.365$.
\textbf{15.} $5.0/0.365 = \qty{13.7}{h}$.
\textbf{16.} 2.7.
\textbf{17.} 6.5, el efecto isotópico completo.
\textbf{18.} Dosis menores o menos frecuentes, y concentraciones más estables en la sangre.
\textbf{19.} $\ln(k/T)$ en función de $1/T$.
\textbf{20.} $\Delta^{\ddagger}H^\circ = 61.9 \pm \qty{0.3}{kJ/mol}$, $\Delta^{\ddagger}S^\circ = -60.1 \pm \qty{0.9}{J.K^{-1}.mol^{-1}}$.
\textbf{21.} $\Delta^{\ddagger}H^\circ = 66.8 \pm \qty{0.2}{kJ/mol}$, $\Delta^{\ddagger}S^\circ = -60.3 \pm \qty{0.7}{J.K^{-1}.mol^{-1}}$.
\textbf{22.} $\Delta\Delta^{\ddagger}H^\circ = 4.88 \pm \qty{0.33}{kJ/mol}$ ($\sqrt{0.26^2 + 0.21^2}$).
\textbf{23.} Sí: 4.83 queda holgadamente dentro de un error típico. El efecto isotópico se
explica por completo por las \omterm{def:b3:quantum-model-systems:oscillator}{energías de punto cero}, sin indicio de \omterm{def:b3:quantum-model-systems:tunnelling}{efecto túnel}.
\textbf{24.} Son iguales dentro de sus errores: el estado de transición tiene la misma
estructura para ambos \omterm{def:b3:rovibrational-spectroscopy:rotational-constant}{isotopólogos}, como se supuso en la Parte I.
\textbf{25.} \textbf{$t_{1/2}(\mathrm D)/t_{1/2}(\mathrm H) \approx 2.7$ a \qty{37}{\celsius}}: \qty{13.7}{h} frente a \qty{5.0}{h}.
\end{solution}
